% ===========================================================================
\section{Related Work}
\label{sec:related}
% ===========================================================================

\subsection{Household energy inventories and their activity problem}

Residential fuel inventories are usually built one of two ways: from a household survey giving
fuel use per household by fuel and region, or from national fuel supply attributed to the
domestic sector. Surveys give richer structure and are expensive and infrequent; supply
statistics are annual and coarse.

Pakistan has both in principle. The Pakistan Social and Living Standards Measurement survey
carries household fuel questions, and OGRA publishes gas and petroleum supply by consuming
sector. This inventory uses the supply route for the fuels that pass through a market, for the
same reason the transport sector does:
the quantity that a national inventory must reproduce is the fuel the country consumed, and a
survey-based estimate has to be reconciled against it afterwards.

The weakness of that choice is specific and is treated in Section~\ref{sec:limitations}: supply
statistics cover fuels that pass through a market, and the two largest household biomass fuels in
Pakistan---firewood and dung---largely do not; for those two the household survey is the estimate.

\subsection{Household air pollution in South Asia}

Solid-fuel combustion in households is among the most studied exposure pathways in the region,
and the health literature on it is far stronger than the emissions literature. What is
comparatively weak is the link between the two: exposure studies measure concentrations in
and around dwellings, while emission inventories estimate mass released, and translating
between them requires assumptions about ventilation, stove type and time-activity patterns
that neither literature settles.

This paper does not attempt that translation. It reports emitted mass on the same activity
basis as the greenhouse gases, which is what an inventory can defensibly supply, and treats
the exposure implication as a ranking of sources rather than a concentration estimate.

\subsection{Emission factors for biomass cookstoves}

The Andreae compilation \citep{andreae2019} and the Akagi compilation \citep{akagi2011}
are the standard field-measurement syntheses for biomass combustion, and the EMEP/EEA
small-combustion chapter \citep{emep2023smallcombustion} is the standard source for
residential fossil fuels. All carry wide ranges, because cookstove emissions depend heavily
on stove design, fuel moisture and operator behaviour, and because measurement campaigns
disagree.

The World Bank's 2024 Pakistan Energy Survey \citep{worldbank2024mtf}, built on a
nationwide 2021--22 Multi-Tier Framework survey, publishes the national stove typology
this literature previously lacked: main-stove shares by locality and province, a
clean-fuel main-stove share of \SI{44.3}{\percent}, rural dominance of traditional and
three-stone/open-fire stoves (\SI{71.8}{\percent}), and stove stacking at
\SI{6.2}{\percent}. It stops short of the fuel-by-device-by-end-use crosswalk a fully
technology-resolved emission inventory needs, which is how this paper uses it: as the
evidence base for a declared technology-mix sensitivity rather than a resolved device
layer.

The biomass particulate and black carbon factors used here come from Andreae's compilation,
the fossil factors from the EMEP/EEA chapter, and the crop-residue rows mirror the
field-burning factors \citep{emep2019} the agriculture sector already uses---one
compilation family per burning layer, but literature values rather than Pakistani
measurement throughout, which is the layer's principal limitation and is stated as such.

\subsection{The biomass \texorpdfstring{\COtwo{}}{CO2} convention}

IPCC reports \COtwo{} from biomass combustion as a memo item outside national totals, on the
reasoning that the carbon was removed from the atmosphere by the plant and its release is part
of a short cycle, with any net loss of biomass stock accounted in the land sector instead.
Methane and nitrous oxide from the same combustion are counted in full, because they are
different molecules with different radiative effects and no corresponding uptake.

The convention is standard and Pakistan's own national submission states it for the same fuels
\citep{pakbur1}. It is worth spelling out here because it produces a result that reads oddly
without it: biomass is \SI{65.4}{\percent} of household energy in this sector and
\SI{24}{\percent} of its reported emissions.
